\documentclass[12pt]{article}

\usepackage{amsmath, amssymb, amsthm}
\usepackage{enumerate}
\usepackage{hyperref}
\usepackage{geometry}
\usepackage{newtxtext}
\usepackage{newtxmath}
\geometry{margin=1in}

\title{Probability Theory - Lecture 3}
\author{}
\date{}

% ----------------------------------------------------------------
\begin{document}
\maketitle

One of the central ideas in analysis and probability is that
``averaging inputs'' and ``applying a nonlinear function''
do not usually commute.

Convex functions bend upward,
and because of this,
applying the function after averaging often gives a smaller answer
than averaging after applying the function.

This phenomenon is captured by \textit{Jensen's inequality}.

In this problem, we will prove the simplest finite version.

% ================================================================
\section*{Questions}

\begin{enumerate}[(a)]

\item Expand:
\[
(x-y)^2.
\]

\item Explain why
\[
(x-y)^2 \geq 0.
\]

\item Rearrange the previous inequality to show that
\[
x^2+y^2 \geq 2xy.
\]

\item Show that
\[
x^2+y^2
\geq
\frac{(x+y)^2}{2}.
\]

(Hint: expand $(x+y)^2$.)

\item Divide both sides by $2$ to deduce that
\[
\frac{x^2+y^2}{2}
\geq
\left(\frac{x+y}{2}\right)^2.
\]

\item Explain in words what the previous inequality says.

\item Let
\[
f(x)=x^2.
\]
Rewrite part (e) in the form
\[
\frac{f(x)+f(y)}{2}
\geq
f\left(\frac{x+y}{2}\right).
\]

\item Let $x_1,x_2,\dots,x_n$ be real numbers.
Their average is
\[
\frac{x_1+x_2+\cdots+x_n}{n}.
\]

Guess the general inequality corresponding to part (g).

\item Show that if
\[
f(x)=x^2,
\]
then
\[
\frac{x_1^2+x_2^2+\cdots+x_n^2}{n}
\geq
\left(
\frac{x_1+x_2+\cdots+x_n}{n}
\right)^2.
\]

(Hint: expand
\[
(x_1+x_2+\cdots+x_n)^2
\]
and compare cross-terms.)

\item A function $f$ can be called \textit{convex} (not a definition) if its graph bends upward.
One version of Jensen's inequality states that for convex functions,
\[
f(\text{average})
\leq
\text{average of } f.
\]

Explain why the function
\[
f(x)=x^2
\]
is convex from the inequalities above.

\end{enumerate}

% ================================================================
\newpage
\section*{Solutions}

Credit: Standard introductory analysis and convexity.

\begin{enumerate}[(a)]

\item Expanding gives
\[
(x-y)^2=x^2-2xy+y^2.
\]

\item Any real number squared is nonnegative, so
\[
(x-y)^2\geq 0.
\]

\item Since
\[
x^2-2xy+y^2\geq 0,
\]
moving the middle term to the other side gives
\[
x^2+y^2\geq 2xy.
\]

\item Expanding,
\[
(x+y)^2=x^2+2xy+y^2.
\]

Using part (c),
\[
2xy\leq x^2+y^2.
\]

Substituting,
\[
(x+y)^2
\leq
x^2+(x^2+y^2)+y^2
=
2x^2+2y^2.
\]

Dividing by $2$ gives
\[
\frac{(x+y)^2}{2}
\leq
x^2+y^2.
\]

\item Dividing both sides by $2$,
\[
\left(\frac{x+y}{2}\right)^2
\leq
\frac{x^2+y^2}{2}.
\]

Equivalently,
\[
\frac{x^2+y^2}{2}
\geq
\left(\frac{x+y}{2}\right)^2.
\]

\item The average of the squares is at least the square of the average.

In other words,
squaring after averaging gives a smaller number than averaging after squaring.

\item Since
\[
f(x)=x^2,
\]
we have
\[
f\left(\frac{x+y}{2}\right)
=
\left(\frac{x+y}{2}\right)^2,
\]
and
\[
\frac{f(x)+f(y)}{2}
=
\frac{x^2+y^2}{2}.
\]

Thus part (e) becomes
\[
\frac{f(x)+f(y)}{2}
\geq
f\left(\frac{x+y}{2}\right).
\]

\item The natural generalization is
\[
\frac{f(x_1)+f(x_2)+\cdots+f(x_n)}{n}
\geq
f\left(
\frac{x_1+x_2+\cdots+x_n}{n}
\right).
\]

This is Jensen's inequality for equal weights.

\item Expanding,
\[
(x_1+\cdots+x_n)^2
=
x_1^2+\cdots+x_n^2
+
2\sum_{i<j}x_ix_j.
\]

Also,
\[
0\leq (x_i-x_j)^2=x_i^2-2x_ix_j+x_j^2,
\]
so
\[
2x_ix_j\leq x_i^2+x_j^2.
\]

Summing over all pairs $i<j$,
every square term appears exactly $(n-1)$ times, so
\[
2\sum_{i<j}x_ix_j
\leq
(n-1)\sum_{i=1}^n x_i^2.
\]

Therefore
\[
(x_1+\cdots+x_n)^2
\leq
n\sum_{i=1}^n x_i^2.
\]

Dividing by $n^2$ gives
\[
\left(
\frac{x_1+\cdots+x_n}{n}
\right)^2
\leq
\frac{x_1^2+\cdots+x_n^2}{n}.
\]

\item The inequalities show that averaging inputs before applying $f(x)=x^2$
always gives a value no larger than averaging the outputs.

Geometrically, this means the graph of $y=x^2$ bends upward.

That upward bending is exactly what convexity means. \qed

\end{enumerate}

\end{document}